Tabular data sharing under privacy constraints is increasingly important for research and collaboration. Synthetic data generators (SDGs) are a promising solution, but synthetic data remains vulnerable to attacks, such as membership inference attacks (MIAs), which aim to determine whether a specific record was part of the training data.
State-of-the-art MIAs are powerful but impractical: they rely on shadow modeling, requiring hundreds of SDG training runs, and need auxiliary data several times larger than the original training set.
Fast proxy metrics like distance to closest record (DCR) are efficient but have limited sensitivity to MIA risk.
We introduce ReMIA (Relative Membership Inference Attack), a practical privacy metric that requires only two SDG training runs and additional data no larger than the original training set.
Rather than predicting whether a record was in the training set, ReMIA generates two synthetic datasets from two source datasets and measures whether a classifier can identify which source a record came from. Experiments across multiple tabular datasets and SDGs show that ReMIA has a sensitivity comparable to state-of-the-art MIAs while being substantially more practical.
We further observe that SDGs can achieve privacy-utility trade-offs that traditional noise-based anonymization methods do not match.
Code is available at \url{https://github.com/aindo-com/remia}.
